\begin{figure}
\centering
\footnotesize
\setlength{\tabcolsep}{4pt}
\begin{tabular}{>{\raggedright\arraybackslash}p{0.20\textwidth}>{\raggedright\arraybackslash}p{0.29\textwidth}>{\raggedright\arraybackslash}p{0.12\textwidth}>{\raggedright\arraybackslash}p{0.30\textwidth}}
\arrayrulecolor{gray!60}\hline
\rowcolor{grammarheadcolor}Construct & Python & Exclusion & Mechanism \\
Non-Boolean test of \pyinline{if} or \pyinline{assert}, conditional expression, comprehension guard, first operand of \pyinline{and} or \pyinline{or}, or operand of \pyinline{not} & Coerced to \pyinline{bool} & Static & Test checked against \pyinline{bool} (\ruleName{if}, \ruleName{assert}, \ruleName{cond}, \ruleName{qual-guard}, \ruleName{and}, \ruleName{or}, \ruleName{unop}) \\\rowgap
Boolean operand of arithmetic or ordering, index of a sequence, or argument of \pyinline{round} & Treated as the integer 0 or 1 & Static & \pyinline{bool} not a subtype of \pyinline{int} (\ruleName{binop}, \ruleName{unop}, \ruleName{subscript-list}, \ruleName{subscript-tuple-int}, \ruleName{subscript-str}, \ruleName{call}) \\\rowgap
Second argument of \pyinline{round} & Rounds to that number of digits & Static & \pyinline{round} takes one argument (\ruleName{call}) \\\rowgap
Argument of \pyinline{sys.exit} other than an integer, or no argument & Exits with status 0 for \pyinline{None} or no argument; otherwise prints the argument and exits with status 1 & Static & Argument checked against \pyinline{int} (\ruleName{call}) \\\rowgap
List pattern at tuple, or tuple pattern at list & Matches, since list and tuple patterns both match lists and tuples & Static & Pattern fails at a type with tuple or list values (\ruleName{check-list-none}, \ruleName{check-tuple-none}) \\\rowgap
List or tuple pattern at \pyinline{range} & Matches, since \pyinline{range} is a sequence & Static & Pattern fails at a type with range values (\ruleName{check-list-none}, \ruleName{check-tuple-none}) \\\rowgap
Positional constructor pattern with fewer sub-patterns than fields & Matches, ignoring remaining fields & Static & Arguments cover fields exactly ($\fieldMap{\Sigma}$) \\\rowgap
Non-string dictionary key & Any hashable value accepted, with \pyinline{1}, \pyinline{1.0} and \pyinline{True} equivalent as keys & Syntactic in patterns, static in expressions & Keys are string literals in patterns, of type \pyinline{str} in expressions (\ruleName{dict}, \ruleName{dict-comp}, \ruleName{subscript-dict} and the \pyinline{in} overload at \pyinline{dict}) \\\rowgap
Non-string \pyinline{assert} message & Any object accepted, shown in traceback & Static & Message checked against \pyinline{str} (\ruleName{assert-msg}) \\\rowgap
\pyinline{==} or \pyinline{!=} between operands of unrelated types & Defined, giving \pyinline{False} & Static & Operand types must be comparable (\ruleName{binop}) \\\rowgap
Boolean against number under \pyinline{==} or \pyinline{!=} & Compared as numbers, so \pyinline{True} equals \pyinline{1} & Static; dynamic at \pyinline{object} & \pyinline{bool} and \pyinline{int} incomparable (\ruleName{binop}); at \pyinline{object}, equality undefined \\\rowgap
Function under \pyinline{==} or \pyinline{!=} & Compared by identity, so two textually identical lambdas are unequal & Static; dynamic at \pyinline{object} & Callable types not equality types (\ruleName{binop}); at \pyinline{object}, undefined between functions \\\rowgap
\pyinline{nan} under container equality or membership & Elements compared by identity before equality, so identical NaNs match but distinct ones do not & Dynamic & Undefined when both elements are \pyinline{nan}
\end{tabular}
\caption{Python coercions and lax matches excluded by \PurePy}
\label{fig:coercions}
\end{figure}
